% ============================================================================
% CP423 Final Project Report
%
% BUILD
%   pdflatex report.tex   (twice)
%
%   No shell escape required. The pipeline diagram is authored in Mermaid
%   (pipeline.mmd) and pre-rendered to pipeline.png. Regenerate with:
%     mmdc -i pipeline.mmd -o pipeline.png -t neutral -b white -s 3
%   Upload report.tex and pipeline.png together on Overleaf.
%
%   TARGET LENGTH: 3 to 4 pages excluding appendices. All figures, tables, and
%   verbatim prompts live in the appendix, which can run as long as needed.
%
%   NOTE ON NUMBERS: every figure in this report matches the tracked files in
%   results/ as of the last full run. That run predates the removal of the
%   automated verifier and the post-generation refusal rewriting described in
%   Section 5. Re-run `python reproduce.py`, re-grade the manual review CSV,
%   and refresh Sections 7 and 8 and Appendices E and F before submitting.
%   Retrieval metrics are unaffected; retrieval did not change.
% ============================================================================

\documentclass[10pt,a4paper]{article}

\usepackage[T1]{fontenc}
\usepackage[utf8]{inputenc}
\usepackage[scaled=0.95]{helvet}
\renewcommand{\familydefault}{\sfdefault}

\usepackage[margin=0.9in]{geometry}
\usepackage{booktabs}
\usepackage{graphicx}
\usepackage{listings}
\usepackage{titlesec}
\usepackage{caption}
\usepackage{enumitem}
\usepackage[hidelinks]{hyperref}

\setlength{\parskip}{0.45em}
\setlength{\parindent}{0pt}
\linespread{1.0}
\setlength{\emergencystretch}{1.5em}
\captionsetup{font=small}

\titleformat{\section}{\normalsize\bfseries}{\thesection}{0.7em}{}
\titleformat{\subsection}{\normalsize\itshape}{\thesubsection}{0.7em}{}
\titlespacing*{\section}{0pt}{1.0em}{0.35em}
\titlespacing*{\subsection}{0pt}{0.7em}{0.25em}

\lstset{
  basicstyle=\ttfamily\footnotesize,
  breaklines=true,
  frame=single,
  framesep=5pt,
  numbers=none,
  showstringspaces=false,
  columns=fullflexible,
  breakatwhitespace=false,
  aboveskip=0.8em,
  belowskip=0.8em
}

% One entry of the evaluation set: number, gold chunks, question, gold answer.
\newcommand{\evalq}[4]{%
  \medskip\noindent{\raggedright#1\ \textit{(#2)}\par
  #3\par
  \textit{Gold answer:} #4\par}}

% ============================================================================

\title{\vspace{-2.5em}\large CP423 Final Project: Retrieval-Augmented Generation over Material UI Documentation}
\author{\normalsize Cameron Grenier, 169069484 \and \normalsize Muhammad Wasay Munawar, 169079709}
\date{\small August 5, 2026}

\begin{document}

\maketitle
\vspace{-1.5em}
\begin{center}
\small
Repository: \url{https://github.com/CameronGrenier/LLM-Search-Engine} \quad\textbar\quad
Demo video: \url{https://VIDEO-LINK}
\end{center}

% ===========================================================================
\section{Introduction and System Overview}
% ===========================================================================

A language model answers documentation questions from a memory frozen at its training cutoff, so when a library ships a new major version the model keeps describing the old API and gives no sign that it is doing so. Retrieval-augmented generation fetches passages from an authoritative corpus at query time and constrains the model to answer from them.

\textit{Research question:} on a documentation corpus that post-dates the generator's training cutoff, does retrieval improve factual accuracy and abstention over a no-context baseline, and how does lexical retrieval compare against dense neural retrieval on identical chunks?

The comparison has three arms: no retrieval, BM25, and dense retrieval. Both indexes are built over the same \texttt{chunks.jsonl}, and the generator, its decoding settings, and its prompt scaffolding are held constant, so any difference in output is attributable to retrieval alone. The pipeline runs in two phases, indexing once over the corpus and query serving per question, and is drawn in Appendix~\ref{app:diagram}.

% ===========================================================================
\section{Corpus Construction}
% ===========================================================================

\subsection{Source and preprocessing}

The corpus is the \href{https://github.com/mui/material-ui}{Material UI documentation at tag v9.2.0}: 151 markdown prose files and 143 API reference JSON files, 294 documents in total. \texttt{fetch\_corpus.py} downloads one codeload tarball rather than walking the GitHub API file by file, which would need several hundred requests and hits unauthenticated rate limits. The tarball is cached and the tag is pinned in \texttt{config.py}, so later upstream edits cannot change the experiment.

MUI was chosen for version drift and verifiability. It ships breaking changes quickly and v9.2.0 post-dates the generator's cutoff, which creates headroom between what the model knows and what the docs say. Every answer we ask for is a prop name, a default value, or a version number checkable against a line of a pinned file, so grading is not a judgment call.

\texttt{preprocess.py} treats the two file types differently. Prose is MDX, so embedded HTML is stripped with BeautifulSoup before anything else runs, since a tag-stripping regex cannot tell an HTML wrapper from a JSX element. Frontmatter supplies title and component metadata. API JSON is flattened into sentences of the form \texttt{Prop `max': type number. Default: 99}, because the retrievers would otherwise match JSON punctuation rather than field meaning, and these files carry most of the factoid answers.

Demo directives get particular handling. MUI embeds interactive examples with a directive naming a source file, and the obvious move is to inline that source at index time. We tried it and rejected it: every demo opens with a near-identical import and export block, so hundreds of chunks came to share a large common component in embedding space, which compresses the distances retrieval depends on, and the added code pushed documents past the embedder's token limit. Each directive is instead replaced by a position marker and the source is stored separately in \texttt{demos.jsonl}, so the index sees prose only while the marker records where the code sat. The code is spliced back into the retrieved passage at generation time, since a question about how to do something in MUI is answered by the example rather than by the prose describing it. Retained per document are title, component, repository path, document type, and a deterministic identifier derived from the filename, which is what makes a citation traceable back to a pinned file.

\subsection{Chunking policy}

Prose is split on headings. Each chunk holds two consecutive heading sections and adjacent chunks overlap by one heading, so no boundary separates a claim from the material framing it. Sections longer than 500 tokens are divided further with a 100-token overlap. API records are kept whole, since a prop name loses meaning once separated from its type, default, and description. This produces 2140 chunks, 1997 prose and 143 API, averaging 104 tokens. Both retrievers index this same set, so chunking is held constant across the comparison. Counts are in Appendix~\ref{app:corpus}.

% ===========================================================================
\section{Corpus Suitability Diagnostic}
% ===========================================================================

RAG is only worth measuring if the corpus holds information the generator does not already have. The ten answerable evaluation questions, 7 factoid and 3 multi-hop, were put to Qwen2.5-3B-Instruct with no retrieved context and the naive system prompt, so nothing suppressed an answer the model could give.

None of the ten was fully correct. Two were partially correct and eight were wrong, and the model volunteered plausible prop names and defaults rather than declining to answer. Adding retrieval improved 9 of the 10: dense RAG answered 8 fully correctly and BM25 answered 3. The corpus is therefore not memorised at v9.2.0, and the gains measured in Section~7 are attributable to retrieval rather than to parametric knowledge. Counts are in Appendix~\ref{app:diagnostic}.

% ===========================================================================
\section{Retrieval}
% ===========================================================================

\subsection{BM25}

BM25 runs over the same 2140 chunks with $k_1 = 1.5$ and $b = 0.75$. Tokenisation is a lowercased word-boundary split and nothing more. There is no stopword list and no lemmatiser, because both retrievers share one chunk text and aggressive normalisation would help the lexical side while damaging the dense side, which needs natural word order.

The two parameters control the two ways BM25 departs from raw TF-IDF. Term frequency saturates under $k_1$, so the tenth occurrence of a term adds almost nothing over the ninth and a chunk cannot win by repetition alone. Length normalisation under $b$ divides by chunk length relative to the corpus average, so a long chunk does not accumulate score merely by containing more words. Both matter here, where API chunks are long and repeat prop names while prose chunks are short.

The index holds a term-frequency counter per chunk plus a corpus-wide document-frequency table, and a query is scored by a linear scan over every chunk. That is naive next to a postings-list lookup, but at this size it answers in milliseconds and needs no model, no GPU, and no re-indexing when the generator changes. BM25 is strongest on exact component and prop names, where a rare token carries high inverse document frequency, and weakest when a question is worded differently from the documentation it needs to find.

\subsection{Dense retrieval}

Dense retrieval embeds chunks and queries into one space and treats relevance as proximity, matching paraphrase where BM25 matches tokens. The model is \texttt{jinaai/jina-embeddings-v3-hf}, 1024 dimensions with an 8192-token limit, chosen for its task-specific LoRA adapters and because it indexes and queries at usable speed on a consumer laptop. Retrieval is asymmetric: corpus text uses the \texttt{retrieval\_passage} adapter and queries use \texttt{retrieval\_query}, projecting both sides into a space where question-shaped text lands near answer-shaped text. Only one adapter may be loaded per process, since loading a second re-initialises the first and LoRA zeroes its B matrix, silently turning the clobbered adapter into a no-op and leaving the two sides in different spaces. The adapter is therefore chosen once through an environment variable that entry-point scripts set before importing the module, and a startup check warns if the loaded weights are all zero.

Chunk vectors come from late chunking. MUI docs contain many near-identical passages, since an install snippet reads the same across fifty component pages, and a chunk embedded in isolation yields a vague vector because its tokens attend only to the chunk. Late chunking instead pools each chunk vector from a forward pass over the surrounding document, so every token has already attended to that context. This wants the whole document in one pass, which the 8192-token limit does not allow for the four longest documents, the largest at 26{,}535 tokens.

Those documents are covered by windows built in two stages. Whole chunks are packed backwards from the document end until coverage would exceed 6000 tokens, so window edges land on chunk boundaries and every chunk is owned by exactly one window. Each window start is then extended backward toward an 8000-token ceiling as read-only context. Deciding ownership at packing time rather than by containment at pooling time matters, because the extended regions overlap and containment would be ambiguous. Both operations run backward because documentation is read top to bottom, so context preceding a chunk carries more signal than context following it.

Vectors go into a FAISS \texttt{IndexFlatIP}, exact brute-force search, since IVF or HNSW would trade recall for speed the system does not need at 2140 vectors. L2 normalisation at pooling time makes inner product equal cosine similarity, and chunk ids are persisted in insertion order because FAISS addresses rows by integer position (Appendix~\ref{app:stats}). The gain is paraphrase matching: dense retrieval put the AvatarGroup \texttt{max} chunk first for a question about limiting how many avatars are shown, where BM25 ranked it sixty-first for want of shared vocabulary, while BM25 wins on rare exact tokens such as \texttt{enterTouchDelay}. Indexing costs 201 seconds and any chunking change forces a full re-embed, against a query cost of one forward pass and a dot product. The 8192-token ceiling is the binding design constraint, and Section~9 returns to it.

% ===========================================================================
\section{Generation}
% ===========================================================================

Generation uses \texttt{Qwen/Qwen2.5-3B-Instruct}, named in \texttt{config.py} and loaded locally through Hugging Face \texttt{transformers} with dtype set to auto. Device selection prefers CUDA, then Apple MPS, then CPU; our runs were on MPS. It was chosen as an ungated checkpoint needing no API key, small enough to run unquantised on a laptop, and instruction-tuned enough to obey a prompt confining it to the supplied passages. Its 32{,}768-token window is not a constraint, since five chunks capped at 500 tokens come to roughly 2500.

Decoding is greedy, \texttt{do\_sample=False} with \texttt{max\_new\_tokens=512}, which is a reproducibility requirement rather than a quality choice, since sampling would make the manual judgments in Section~7 unrepeatable. \texttt{src/seeding.py} applies seed 42 from \texttt{config.py} to the Python, NumPy, and Torch generators when either model loads, and \texttt{reproduce.py} sets \texttt{PYTHONHASHSEED} in every stage subprocess, the only place it can take effect. No stage is actually stochastic today, so determinism comes from greedy decoding and the seeding is insurance against that changing.

\texttt{TOP\_K} is 5. Retrieved chunks are pasted in as numbered passages of the form \texttt{[n] (chunk\_id)}, followed by the chunk text and then any demo source that chunk referenced, as fenced code blocks. Appending the code inside the passage rather than as a separate block keeps one citation pointing at one source. Roughly a third of chunks carry a demo, so this raises the mean prompt from about 900 tokens to about 2400, with the largest observed at 5414, all comfortably inside the model's window. Model, device, decoding settings, seed, passage formatting, and system prompt are identical across the BM25 and dense conditions, and only the retrieved set differs. The naive condition needs a different system prompt, because the RAG prompt confines the model to supplied passages and would otherwise abstain for reasons unrelated to its knowledge, invalidating the Section~3 diagnostic. Both prompts are in Appendix~\ref{app:prompts}.

One design decision deserves stating plainly, because it makes our numbers look worse. An earlier version post-processed the answer, replacing it with \texttt{I don't know} whenever it carried no valid citation or hedged, and ran a second pass in which the same model verified its own answer against its own citations. Both were removed. The rewriting made the inline citation rate true by construction rather than measured, since an uncited answer could not reach the output at all, and the verifier approved a reply consisting of nothing but the string \texttt{[1],[2]} as supported by the passages it cited. A 3-billion-parameter model grading its own work shares its own blind spots. The refusal instruction now lives in the system prompt, the model's reply is recorded verbatim, and grounding is judged by a human.

% ===========================================================================
\section{Evaluation Set}
% ===========================================================================

The gold set is 13 questions written by hand while reading the corpus: 7 factoid, 3 multi-hop, and 3 unanswerable. Each carries a reference answer, a type, and the ground-truth chunk ids that support it. Multi-hop questions require evidence from two chunks, so a system must retrieve both to answer. Unanswerable questions have no valid supporting chunk anywhere in the corpus and take \texttt{I don't know} as the reference answer, testing abstention rather than recall. The full set is in Appendix~\ref{app:evalset}.

Several sources of bias should be read alongside the results. The set is small, so one question shifts any metric by roughly eight percentage points. Component coverage is uneven and reflects what was interesting to read rather than a sample of the library. The questions were written and graded by the same people who built the system, and scoring is manual and therefore subjective. Exact gold matching also penalises a retriever that returns a neighbouring chunk carrying comparable evidence, which our one-heading overlap makes common.

% ===========================================================================
\section{Results}
% ===========================================================================

Retrieval was scored on the ten answerable questions. Dense retrieval beat BM25 on all four metrics, most clearly on recall at 0.800 against 0.650 and on nDCG@5 at 0.669 against 0.584. MRR was nearly tied at 0.683 against 0.674, which fits the design: both systems find the same easy chunks near the top, and dense retrieval finds more of the remaining evidence further down the ranking.

All 39 answers were scored manually from 0 to 2 on correctness, groundedness, and completeness. Dense RAG was fully correct on 10 of 13 questions, BM25 on 6, and the no-context baseline on none. Dense led on every dimension, at 1.69 for correctness and completeness against 1.15, and 1.77 for groundedness against 1.69. The groundedness gap is much the narrower of the two, which is consistent with how the retrievers work: BM25 returns chunks containing the query's exact wording, so a citation lands on text that visibly supports the claim, while dense retrieval returns paraphrase matches whose support is real but less literal. Full tables are in Appendix~\ref{app:results}.

The unanswerable questions separated the two systems. BM25 RAG declined all three using the exact fallback wording. Dense RAG declined two, and on the Card question answered that the passages did not define the prop rather than using the required string, which is the correct judgment expressed in a form the exact-match metric scores as a non-refusal. The baseline declined none and fabricated plausible detail for each. Retrieval therefore changes not only what the system knows but whether it recognises the boundary of what it knows.

% ===========================================================================
\section{Error Analysis}
% ===========================================================================

Failures separate into three kinds with different causes. Retrieval failures mean the evidence never reached the model, generation failures mean it reached the model and was not used, and citation failures mean the answer was right but not traceable to the passage supporting it.

\textit{Context present but unused.} On Q1 the gold chunk ranked first under both retrievers and contains the install command verbatim, and both systems still answered that Base UI must be installed without naming the package. On Q2 under BM25 the gold chunk again ranked first and the model's entire reply was \texttt{[1],[2]}, citations with no content. On Q6 the Badge API chunk carrying the default of 99 ranked second under dense retrieval, but the model built its answer from a Rating API chunk two places below and reported Rating's default of 5 instead. The last is the most instructive, since two flattened prop tables are structurally near-identical and the prompt rule against inferring defaults from related components did not prevent the substitution. The evidence reached the prompt in all three cases, so no change to chunking, embedding, or ranking would fix them.

\textit{Multi-hop retrieval.} This was the dominant weakness. Five of the six question-by-system pairs recovered only one of the two required chunks, and the pattern is consistent: the first hop ranks near the top and the second sits far down. On the Rating question BM25 ranked the prose hop first and the API hop twenty-third and declined to answer, while dense retrieved both and answered correctly. On the Slider question BM25 ranked the prose hop first and the API hop one hundred and eleventh, so it named \texttt{shiftStep} correctly and could not supply the defaults. The Masonry question failed for both, with the ImageList evidence at rank eight under BM25 and twenty-one under dense. The cause is structural, since both retrievers rank a flat list against a single query representation, so a two-hop question is dominated by its first hop and the second competes against chunks matching the first. Raising $k$ surfaces the missing evidence but dilutes the prompt. Dense on the Slider question shows the limit of the metric from the other side: recall@5 was 0.5 and the answer was still fully correct, because the one API chunk it retrieved carried both defaults.

\textit{Citation faithfulness.} Citations frequently land on a neighbouring chunk rather than the labelled gold one, with mean citation precision at 0.50 for BM25 and 0.55 for dense, and recall at 0.35 and 0.55. On Q5 the gold chunk ranked first under both retrievers and both cited the overlapping chunk after it. Those answers are correct and the cited text does support them, so the gap is partly real and partly our gold labels being stricter than our own chunking. Separately, every citation record initially carried a null \texttt{doc\_id}, because the citation builder read the field defensively while neither retriever emitted it, reducing a citation to a bare chunk id on all 39 records. Both retrievers now return \texttt{doc\_id} and the field is read directly, so the same mismatch would fail immediately rather than write nulls into every record.

% ===========================================================================
\section{Limitations and Future Work}
% ===========================================================================

The binding constraint on the dense side is jina-v3's 8192-token window. Four documents exceed it, the largest at 26{,}535 tokens, and they require the backward-packing windows of Section~4.2, which complicate indexing and leave those chunks without whole-document context. A larger-context embedder such as jina-embeddings-v4 or Qwen3-Embedding-8B, both around 32K tokens, would let every document in this corpus embed in one pass. The cost is real: both are substantially larger models, so memory and indexing time rise, adopting one requires a full re-embed, and a 32K window moves the wall rather than removing it.

The multi-hop failures point at concrete work rather than a redesign. Hybrid retrieval fusing the BM25 and dense rankings would recover evidence that one retriever ranks well and the other does not, which fits the Masonry case where BM25 ranked the missing chunk thirteen places higher. Reranking a larger candidate pool with a cross-encoder would let the system look past the top five without paying the prompt-dilution cost. Retrieving neighbouring chunks alongside each hit, or decomposing a multi-hop question into subqueries, attacks the same problem from the query side.

The rest are limits of scope: a single technical domain, a 13-question evaluation set, a fixed top-five cutoff, manual and subjective scoring, and the capacity of a 3-billion-parameter generator, which Section~8 shows directly in the cases where correct evidence reached the prompt and went unused.

% ===========================================================================
\section{Conclusion}
% ===========================================================================

Retrieval substantially improved answer quality on this corpus. Dense retrieval beat BM25 on every retrieval metric, and dense RAG was fully correct on 10 of 13 questions against 6 for BM25 and none without context. The diagnostic showed the same effect from the other direction, with retrieval fixing 9 of the 10 answerable questions the generator could not answer from memory, and both RAG systems recognised the unanswerable questions where the baseline fabricated an answer to each.

The result is not that dense retrieval dominates. BM25 remains cheaper by every practical measure, stronger on exact API names, more reliable at producing the required refusal wording, and it ranked the missing Masonry evidence well above dense retrieval. Both systems depended on retrieving every required chunk within the top five, which is where multi-hop questions broke and where hybrid retrieval or reranking would help most.

% ===========================================================================
\appendix
\clearpage
\section{System Pipeline Diagram}
\label{app:diagram}
% ===========================================================================

\begin{figure}[h]
  \centering
  \includegraphics[width=\linewidth,height=0.7\textheight,keepaspectratio]{pipeline.png}
  \caption{System pipeline. Indexing runs once over the corpus. Query serving branches on retrieval condition and converges on a single generator. Authored in Mermaid; source in \texttt{pipeline.mmd}.}
  \label{fig:pipeline}
\end{figure}

\clearpage

% ===========================================================================
\section{Prompts}
\label{app:prompts}
% ===========================================================================

System prompt, RAG conditions:

\begin{lstlisting}
Answer the question using only the supplied numbered passages.

If the passages do not explicitly answer the question, your entire reply must be
this exact sentence:

I don't know

Write it exactly that way: a straight apostrophe, no trailing period, no citation,
no explanation, and nothing before or after it.

Otherwise, answer the question and follow these rules:
- Answer every part of the question that is explicitly supported.
- Place an inline citation such as [1] immediately after every factual sentence.
- Use only passage numbers that were supplied.
- Do not use prior knowledge.
- Do not infer prop names, defaults, variants, or component behaviour from related components.
- Do not treat a passing mention, dependency list, or example name as an answer.
- Never reply with citation markers alone; every reply must contain a sentence.
- Keep the answer concise.
\end{lstlisting}

System prompt, naive condition:

\begin{lstlisting}
Answer the question about Material UI as accurately as you can.
Keep it concise and to the point.
\end{lstlisting}

User message, RAG conditions. Each passage is the chunk text prefixed by its
number and chunk id, followed by any demo source that chunk referenced:

\begin{lstlisting}
Numbered passages:
[1] (prose__accordion__chunk_0004)
<chunk text>

```jsx
<demo source referenced by this chunk>
```

[2] (api__tooltip__chunk_0001)
<chunk text>
...

Question:
<question text>

Answer:
\end{lstlisting}

% ===========================================================================
\section{Corpus and Chunk Statistics}
\label{app:corpus}
% ===========================================================================

\begin{table}[h]
\centering
\small
\begin{tabular}{lr}
\toprule
Source type & Count \\
\midrule
Markdown files & 151 \\
API JSON files & 143 \\
Referenced demo files & 1{,}567 \\
\midrule
Searchable documents & 294 \\
\bottomrule
\end{tabular}
\hspace{2em}
\begin{tabular}{lr}
\toprule
Chunk type & Count \\
\midrule
Prose & 1{,}997 \\
API & 143 \\
\midrule
Total & 2{,}140 \\
\bottomrule
\end{tabular}
\caption{Corpus sources and resulting chunk counts. Average chunk length is 104 whitespace tokens; 613 of the extracted demo files are referenced by a directive and stored in \texttt{demos.jsonl}.}
\label{tab:corpus}
\end{table}

% ===========================================================================
\section{Corpus Suitability Diagnostic Results}
\label{app:diagnostic}
% ===========================================================================

\begin{table}[h]
\centering
\small
\begin{tabular}{lr}
\toprule
No-context result & Count \\
\midrule
Fully correct & 0 \\
Partially correct & 2 \\
Incorrect & 8 \\
\bottomrule
\end{tabular}
\hspace{2em}
\begin{tabular}{lr}
\toprule
System & Fully correct \\
\midrule
No context & 0/10 \\
BM25 RAG & 3/10 \\
Dense RAG & 8/10 \\
\bottomrule
\end{tabular}
\caption{Diagnostic on the ten answerable questions, 7 factoid and 3 multi-hop. Left: breakdown with no retrieved context. Right: fully correct answers by condition.}
\label{tab:diagnostic}
\end{table}

% ===========================================================================
\section{Evaluation Set}
\label{app:evalset}
% ===========================================================================

The 13 hand-written questions with their gold answers and the chunks that
support them. Multi-hop questions require evidence from two chunks.
Unanswerable questions have no supporting chunk anywhere in the corpus.

\subsection*{Factoid}

\evalq{Q1}{prose\_\_menubar}
  {Does Material UI include a Menubar component out of the box, and what do I need to install to use the menubar shown in the docs?}
  {No. Material UI does not ship a menu bar component. The docs page composes the Base UI Menubar styled to Material Design, so Base UI must be installed first with \texttt{npm install @base-ui/react}.}

\evalq{Q2}{prose\_\_progress}
  {Which prop keeps the CircularProgress track visible at all times?}
  {The \texttt{enableTrackSlot} prop.}

\evalq{Q3}{prose\_\_buttons}
  {From which Material UI version is the loading prop available on Button?}
  {Starting from v6.4.0.}

\evalq{Q4}{prose\_\_tooltips}
  {Why does showing a tooltip on touch devices require a long press?}
  {Because \texttt{enterTouchDelay} defaults to 700\,ms.}

\evalq{Q5}{prose\_\_avatars}
  {How do I customize the surplus avatar in an AvatarGroup, and what argument does the callback receive?}
  {Pass a callback to \texttt{renderSurplus}. It receives the surplus number, derived from the children and the \texttt{max} prop, and returns a \texttt{React.ReactNode}.}

\evalq{Q6}{api\_\_badge}
  {Which Badge prop caps large numeric values, and what value does it cap at by default?}
  {The \texttt{max} prop, which defaults to 99.}

\evalq{Q7}{api\_\_avatar-group}
  {Which prop limits how many avatars an AvatarGroup shows, and what is its default value?}
  {The \texttt{max} prop, which defaults to 5.}

\subsection*{Multi-hop}

\evalq{Q8}{prose\_\_rating, api\_\_rating}
  {Which Rating prop sets the minimum increment the value can change by, and what is that prop's default?}
  {The \texttt{precision} prop, which defaults to 1.}

\evalq{Q9}{prose\_\_slider, api\_\_slider}
  {Which Slider prop controls stepping with Page Up/Down or Shift plus arrow keys, and what is its default compared to the default step?}
  {The \texttt{shiftStep} prop, which defaults to 10 while \texttt{step} defaults to 1, so the default \texttt{shiftStep} is ten times larger.}

\evalq{Q10}{prose\_\_masonry, prose\_\_image-list}
  {Masonry orders its children by row. Which component should I use instead to order images by column, and what determines the container heights there?}
  {Use ImageList, specifically the masonry image list variant, where container heights are sized dynamically to reflect each image's aspect ratio.}

\subsection*{Unanswerable}

\evalq{Q11}{none}
  {What does the type prop do in the Card MUI component?}
  {I don't know.}

\evalq{Q12}{none}
  {How do I use the TreeView MUI component?}
  {I don't know.}

\evalq{Q13}{none}
  {The MUI Button component comes with four variants: text, contained, outlined, and one other. What is this fourth variant?}
  {I don't know.}

% ===========================================================================
\section{Full Results Tables}
\label{app:results}
% ===========================================================================

\begin{table}[h]
\centering
\small
\begin{tabular}{lrrrr}
\toprule
Retriever & Precision@5 & Recall@5 & MRR & nDCG@5 \\
\midrule
BM25 & 0.1600 & 0.6500 & 0.6736 & 0.5839 \\
Dense & 0.2000 & 0.8000 & 0.6827 & 0.6693 \\
\bottomrule
\end{tabular}
\caption{Retrieval metrics on the ten answerable questions. The three unanswerable questions are excluded, since they have no gold chunk.}
\label{tab:retrieval}
\end{table}

\begin{table}[h]
\centering
\small
\begin{tabular}{lrrrr}
\toprule
System & Correctness & Groundedness & Completeness & Fully correct \\
\midrule
No context & 0.1538 & 0.1538 & 0.1538 & 0/13 \\
BM25 RAG & 1.1538 & 1.6923 & 1.1538 & 6/13 \\
Dense RAG & 1.6923 & 1.7692 & 1.6923 & 10/13 \\
\bottomrule
\end{tabular}
\caption{Manual scores from 0 to 2 across all 13 questions per condition.}
\label{tab:generation}
\end{table}

\begin{table}[h]
\centering
\small
\begin{tabular}{lrrr}
\toprule
Metric & No context & BM25 RAG & Dense RAG \\
\midrule
Exact match & 0.0000 & 0.2308 & 0.1538 \\
Mean token F1 & 0.4250 & 0.5277 & 0.6252 \\
Unanswerable refusal accuracy & 0.000 & 1.000 & 0.667 \\
Mean citation precision & n/a & 0.500 & 0.545 \\
Mean citation recall & n/a & 0.350 & 0.550 \\
\bottomrule
\end{tabular}
\caption{Automated generation metrics. Refusal accuracy requires the exact fallback string, which is why dense retrieval scores 0.667 despite declining to answer all three unanswerable questions in substance.}
\label{tab:autogen}
\end{table}

% ===========================================================================
\section{Dense Index Statistics}
\label{app:stats}
% ===========================================================================

\begin{table}[h]
\centering
\small
\begin{tabular}{lr}
\toprule
Metric & Value \\
\midrule
Documents & 294 \\
Chunks & 2140 \\
Documents embedded in a single pass & 290 \\
Documents requiring windows & 4 \\
Forward passes over windows & 305 \\
Chunks pooled from a window & 2140 \\
Chunks embedded independently (fallback) & 0 \\
Embedding dimension & 1024 \\
Longest document & 26{,}535 tokens \\
Indexing time & 201\,s (Apple Silicon, MPS) \\
\bottomrule
\end{tabular}
\caption{Dense index statistics from the final clean run.}
\label{tab:densestats}
\end{table}

\end{document}
